\documentclass[10pt]{extarticle}
\usepackage[a4paper,margin=1.5cm]{geometry}
\usepackage[T1]{fontenc}
\usepackage{amsmath,amssymb}
\usepackage{graphicx}
\usepackage[export]{adjustbox}
\usepackage{multicol}
\usepackage{enumitem}
\usepackage{booktabs,array,tabularx}
\setlength{\columnsep}{0.7cm}
\setlength{\parindent}{0pt}
\setlength{\parskip}{3pt}
\pagestyle{empty}
\makeatletter
% figure registry (figdefs.tex is generated by the extract step)
\newcommand{\qfigdef}[4]{\@namedef{qf@file@#1}{#2}\@namedef{qf@w@#1}{#3}\@namedef{qf@wide@#1}{#4}}
\newcommand{\qf@check}[1]{\@ifundefined{qf@file@#1}{\PackageError{score2}{Unknown figure id #1}{}}{}}
\newcommand{\qf@img}[4][]{\includegraphics[width=#2\dimexpr\@nameuse{qf@w@#4}pt\relax,max width=#3,max height=0.45\textheight,#1]{\@nameuse{qf@file@#4}}}
% Full-width questions: finalize wraps every question that contains a wide figure in \qfullbegin...\qfullend,
% so the WHOLE question (stem, figures, options) is set across the page in reading order, then the two columns
% resume. The question's head - its number, stem and first figure (\qfullsplit marks the end of the head) -
% is kept on one page, so a number is never left at a page foot with its figure overleaf; the rest may break
% across pages normally, so no big gaps are left. Without a \qfullsplit, the whole question is kept together
% if it is under 45% of a page. Never write these yourself.
\newif\ifqf@full
\newif\ifqf@boxed
\newsavebox{\qf@box}
\newcommand{\qf@open}{\global\qf@fulltrue\par\global\qf@boxedtrue
  \setbox\qf@box\vbox\bgroup\hsize\textwidth\linewidth\textwidth}
\newcommand{\qf@emit}[1]{\par\egroup\global\qf@boxedfalse
  \ifdim\dimexpr\ht\qf@box+\dp\qf@box\relax<#1\textheight\noindent\usebox{\qf@box}\par\else\unvbox\qf@box\fi}
\newcommand{\qfullsplit}{\ifqf@boxed\qf@emit{0.85}\fi}
\newcommand{\qf@close}{\ifqf@boxed\qf@emit{0.45}\fi\global\qf@fullfalse}
\newcommand{\qfullbegin}{\end{multicols}\qf@open}
\newcommand{\qfullend}{\qf@close\begin{multicols}{2}}
\newcommand{\qfullstart}{\qf@open}
\newcommand{\qfulldone}{\qf@close}
% \qfig[scale]{id}: block figure, centred, directly where it is written (size set by the tool).
% Fallback if a wide figure is not inside a full-width question: break out of the columns in place - no floats.
\newcommand{\qfig}[2][1]{\qf@check{#2}%
  \ifqf@full\qf@blk{#1}{#2}\else
  \ifnum\@nameuse{qf@wide@#2}=1 \expandafter\@firstoftwo\else\expandafter\@secondoftwo\fi
  {\end{multicols}\begin{center}\qf@img{#1}{\textwidth}{#2}\end{center}\begin{multicols}{2}}%
  {\qf@blk{#1}{#2}}\fi}
\newcommand{\qf@blk}[2]{\par\smallskip{\centering\qf@img{#1}{\linewidth}{#2}\par}\smallskip}
% \qfigi[scale]{id}: inline figure (inside option lists / tables / minipages); never breaks columns.
\newcommand{\qfigi}[2][1]{\qf@check{#2}\qf@img[valign=c]{#1}{\linewidth}{#2}}
\makeatother
% structure macros - lint reads these to check completeness and figure placement
\newcommand{\papersec}[1]{\par\bigskip{\centering\Large\bfseries PAPER-#1\par}\medskip}
\newcommand{\partsec}[2]{\par\medskip\noindent{\large\bfseries PART-#1: #2}\par\smallskip}
\newcommand{\sect}[1]{\par\smallskip\noindent\textbf{#1}\par}
\newcommand{\qn}[1]{\par\smallskip\noindent\textbf{#1.}~}
\newcommand{\paraq}[2]{\par\smallskip\noindent\textbf{Paragraph for Questions #1 and #2}\par}
\newenvironment{qoptions}{\begin{itemize}[leftmargin=1.8em,itemsep=1pt,topsep=2pt,label={}]}{\end{itemize}}
\newcommand{\opt}[2]{\item[(#1)] #2}
% figure registry (generated by score2_tools.py finalize)
\qfigdef{p1_phy_s1_q04_1}{Star Review (Other) - Test 02 (10 Oct 2022)-img/p1_phy_s1_q04_1.png}{151.4}{0}
\qfigdef{p1_phy_s2_q07_1}{Star Review (Other) - Test 02 (10 Oct 2022)-img/p1_phy_s2_q07_1.png}{106.8}{0}
\qfigdef{p1_phy_s2_q08_1}{Star Review (Other) - Test 02 (10 Oct 2022)-img/p1_phy_s2_q08_1.png}{76.9}{0}
\qfigdef{p1_phy_s2_q08_2}{Star Review (Other) - Test 02 (10 Oct 2022)-img/p1_phy_s2_q08_2.png}{72.7}{0}
\qfigdef{p1_phy_s2_q08_3}{Star Review (Other) - Test 02 (10 Oct 2022)-img/p1_phy_s2_q08_3.png}{90.2}{0}
\qfigdef{p1_phy_s2_q08_4}{Star Review (Other) - Test 02 (10 Oct 2022)-img/p1_phy_s2_q08_4.png}{63.2}{0}
\qfigdef{p1_phy_s2_q09_1}{Star Review (Other) - Test 02 (10 Oct 2022)-img/p1_phy_s2_q09_1.png}{92.3}{0}
\qfigdef{p1_phy_s2_q10_1}{Star Review (Other) - Test 02 (10 Oct 2022)-img/p1_phy_s2_q10_1.png}{108.0}{0}
\qfigdef{p1_phy_s2_q10_2}{Star Review (Other) - Test 02 (10 Oct 2022)-img/p1_phy_s2_q10_2.png}{81.8}{0}
\qfigdef{p1_phy_s2_q10_3}{Star Review (Other) - Test 02 (10 Oct 2022)-img/p1_phy_s2_q10_3.png}{82.0}{0}
\qfigdef{p1_phy_s2_q10_4}{Star Review (Other) - Test 02 (10 Oct 2022)-img/p1_phy_s2_q10_4.png}{67.1}{0}
\qfigdef{p1_phy_s3_q02_1}{Star Review (Other) - Test 02 (10 Oct 2022)-img/p1_phy_s3_q02_1.png}{92.3}{0}
\qfigdef{p1_phy_s3_q03_1}{Star Review (Other) - Test 02 (10 Oct 2022)-img/p1_phy_s3_q03_1.png}{121.8}{0}
\qfigdef{p1_phy_s3_q04_1}{Star Review (Other) - Test 02 (10 Oct 2022)-img/p1_phy_s3_q04_1.png}{92.1}{0}
\qfigdef{p1_phy_s3_q05_1}{Star Review (Other) - Test 02 (10 Oct 2022)-img/p1_phy_s3_q05_1.png}{105.3}{0}
\qfigdef{p1_chem_s1_q02_A}{Star Review (Other) - Test 02 (10 Oct 2022)-img/p1_chem_s1_q02_A.png}{57.1}{0}
\qfigdef{p1_chem_s1_q02_B}{Star Review (Other) - Test 02 (10 Oct 2022)-img/p1_chem_s1_q02_B.png}{74.2}{0}
\qfigdef{p1_chem_s1_q02_C}{Star Review (Other) - Test 02 (10 Oct 2022)-img/p1_chem_s1_q02_C.png}{83.8}{0}
\qfigdef{p1_chem_s3_q05_1}{Star Review (Other) - Test 02 (10 Oct 2022)-img/p1_chem_s3_q05_1.png}{66.1}{0}
\qfigdef{p1_math_s2_q10_1}{Star Review (Other) - Test 02 (10 Oct 2022)-img/p1_math_s2_q10_1.png}{75.1}{0}
\qfigdef{p2_phy_s1_q02_1}{Star Review (Other) - Test 02 (10 Oct 2022)-img/p2_phy_s1_q02_1.png}{37.7}{0}
\qfigdef{p2_phy_s1_q03_1}{Star Review (Other) - Test 02 (10 Oct 2022)-img/p2_phy_s1_q03_1.png}{122.2}{0}
\qfigdef{p2_phy_s2_q05_1}{Star Review (Other) - Test 02 (10 Oct 2022)-img/p2_phy_s2_q05_1.png}{115.8}{0}
\qfigdef{p2_phy_s2_q06_1}{Star Review (Other) - Test 02 (10 Oct 2022)-img/p2_phy_s2_q06_1.png}{62.9}{0}
\qfigdef{p2_phy_s2_q07_1}{Star Review (Other) - Test 02 (10 Oct 2022)-img/p2_phy_s2_q07_1.png}{102.7}{0}
\qfigdef{p2_phy_s2_q08_1}{Star Review (Other) - Test 02 (10 Oct 2022)-img/p2_phy_s2_q08_1.png}{82.0}{0}
\qfigdef{p2_phy_s2_q09_1}{Star Review (Other) - Test 02 (10 Oct 2022)-img/p2_phy_s2_q09_1.png}{195.9}{0}
\qfigdef{p2_phy_s3_q02_1}{Star Review (Other) - Test 02 (10 Oct 2022)-img/p2_phy_s3_q02_1.png}{134.7}{0}
\qfigdef{p2_phy_s3_q08_1}{Star Review (Other) - Test 02 (10 Oct 2022)-img/p2_phy_s3_q08_1.png}{97.7}{0}
\qfigdef{p2_chem_s1_q01_1}{Star Review (Other) - Test 02 (10 Oct 2022)-img/p2_chem_s1_q01_1.png}{39.6}{0}
\qfigdef{p2_chem_s2_q07_3}{Star Review (Other) - Test 02 (10 Oct 2022)-img/p2_chem_s2_q07_3.png}{38.7}{0}
\qfigdef{p2_chem_s2_q07_2}{Star Review (Other) - Test 02 (10 Oct 2022)-img/p2_chem_s2_q07_2.png}{43.5}{0}
\qfigdef{p2_chem_s2_q07_1}{Star Review (Other) - Test 02 (10 Oct 2022)-img/p2_chem_s2_q07_1.png}{47.8}{0}
\qfigdef{p2_chem_s3_q08_1}{Star Review (Other) - Test 02 (10 Oct 2022)-img/p2_chem_s3_q08_1.png}{232.9}{0}
\raggedcolumns
\makeatletter\@beginparpenalty=10000\makeatother
\newcolumntype{L}[1]{>{\raggedright\arraybackslash}p{#1}}
\begin{document}
\begin{center}{\Large\bfseries Star Review (Other) --- Test 02 (10 Oct 2022)}\end{center}
\begin{multicols}{2}
\raggedcolumns
\papersec{1}
\partsec{1}{PHYSICS}
\sect{SECTION-I(i)}
\qn{1} Select the incorrect alternative(s) :-
\begin{qoptions}
\opt{A}{Work done by kinetic friction on a body always results in a loss of its kinetic energy.}
\opt{B}{Work done on a body, in the motion of that body over a close loop is zero for every force in nature.}
\opt{C}{Total mechanical energy of a system is always conserved no matter what type of internal and external forces on the body are present.}
\opt{D}{When total work done by a conservative force on the system is positive then the potential energy associated with this force decreases.}
\end{qoptions}
\qn{2} A particle is moving in the vertical plane. It is attached at one end of a string of length $\ell$ whose other end is fixed. The velocity at the lowest point is $u$. The tension in the string is $\vec T$ and velocity of the particle is $\vec v$ at any position. Then, which of the following quantity will not remain constant ?
\begin{qoptions}
\opt{A}{$\vec T\cdot\vec v$}
\opt{B}{kinetic energy}
\opt{C}{Gravitational potential energy}
\opt{D}{$\vec T\times\vec v$}
\end{qoptions}
\qn{3} For a two-body system in absence of external forces, the kinetic energy as measured from ground frame is $K_O$ and from center of mass frame is $K_{cm}$. Pick up the correct statement(s) :-
\begin{qoptions}
\opt{A}{The kinetic energy as measured from center of mass frame is least.}
\opt{B}{Only the portion of energy $K_{cm}$ can be transformed from one form to another due to internal changes in the system.}
\opt{C}{The system always retains at least $K_O-K_{cm}$ amount of kinetic energy as measured from ground frame irrespective of any kind of internal changes in the system.}
\opt{D}{The system always retains at least $K_{cm}$ amount of kinetic energy as measured from ground frame irrespective of any kind of internal changes in the system.}
\end{qoptions}
\qn{4} Mass A hits B inelastically ($e=0$) and sticks while moving horizontally with some velocity along the common line of centres of the three equal masses each of same mass. Initially mass B and C are stationary and the spring is unstretched. Then which is correct ?
\qfig{p1_phy_s1_q04_1}
\begin{qoptions}
\opt{A}{Compression will be maximum when blocks have same velocity.}
\opt{B}{Velocity of C is maximum when $(A+B)$ is at rest.}
\opt{C}{Velocity of C is maximum when spring is underformed.}
\opt{D}{Velocity of C is minimum when spring is underformed.}
\end{qoptions}
\qn{5} A wagon filled with sand has a hole so that sand leaks through the bottom at a constant rate $\lambda$. An external force $\vec F$ acts on the wagon in the direction of motion. Assuming instantaneous velocity of the wagon to be $\vec v$ and initial mass of system to be $m_0$, select the equations which does not described the motion of wagon.
\begin{qoptions}
\opt{A}{$\vec F=m_0\dfrac{d\vec v}{dt}+\lambda\vec v$}
\opt{B}{$\vec F=m_0\dfrac{d\vec v}{dt}-\lambda\vec v$}
\opt{C}{$\vec F=(m_0-\lambda t)\dfrac{d\vec v}{dt}$}
\opt{D}{$\vec F=(m_0-\lambda t)\dfrac{d\vec v}{dt}+\lambda\vec v$}
\end{qoptions}
\qn{6} Choose the incorrect statement(s) :
\begin{qoptions}
\opt{A}{Perpendicular axis theorem can be applied for sphere.}
\opt{B}{Perpendicular axis theorem can be applied for cylinder.}
\opt{C}{Perpendicular axis theorem can be applied for disc.}
\opt{D}{Perpendicular axis theorem can be applied for cone.}
\end{qoptions}
\sect{SECTION-I(ii)}
\qn{7} Consider two spherical masses A and B having same density as shown in figure. Spheres are released from rest in free space at an initial separation 21 m. The only interaction force acting on two spheres due to each other is constant in magnitude of 2 N and attractive in nature. Collision is inelastic with coefficient of restitution $\left(\dfrac12\right)$. Mass of B is 1 kg. If magnitude of velocities of A, B before and after the collision are $v_1$, $v_2$ and $v_1'$, $v_2'$.
\qfig{p1_phy_s2_q07_1}
{\small\setlength{\tabcolsep}{3pt}\begin{tabular}{|c|c|c|c|}
\hline
 & \textbf{List--I} & & \textbf{List--II}\\ \hline
(I) & $v_1$ & (P) & $\dfrac12$ m/s\\ \hline
(II) & $v_2$ & (Q) & 8 m/s\\ \hline
(III) & $v_1'$ & (R) & 4 m/s\\ \hline
(IV) & $v_2'$ & (S) & 2 m/s\\ \hline
 & & (T) & 1 m/s\\ \hline
\end{tabular}}\par
\begin{qoptions}
\opt{A}{$\mathrm I\to T;\mathrm{II}\to Q;\mathrm{III}\to P;\mathrm{IV}\to R$}
\opt{B}{$\mathrm I\to R;\mathrm{II}\to Q;\mathrm{III}\to S;\mathrm{IV}\to P$}
\opt{C}{$\mathrm I\to T;\mathrm{II}\to R;\mathrm{III}\to S;\mathrm{IV}\to P$}
\opt{D}{$\mathrm I\to S;\mathrm{II}\to Q;\mathrm{III}\to R;\mathrm{IV}\to T$}
\end{qoptions}
\qn{8} List-I contains some continuous mass system and list-II contain their moment of inertia about axis shown in list-I. Match the following list.\par
{\footnotesize\setlength{\tabcolsep}{2pt}\begin{tabular}{|c|p{0.46\linewidth}|c|c|}
\hline
 & \textbf{List--I} & & \textbf{List--II}\\ \hline
(I) & \qfigi{p1_phy_s2_q08_1} Uniform rod & (P) & $\dfrac{8MR^2}{11}$\\ \hline
(II) & \qfigi{p1_phy_s2_q08_2} Uniform semicircular ring. Axis is perpendicular to plane of ring. & (Q) & $\dfrac{MR^2}{12}$\\ \hline
(III) & \qfigi{p1_phy_s2_q08_3} Uniform Triangular plate of mass $M$ & (R) & $\dfrac{13MR^2}8$\\ \hline
(IV) & \qfigi{p1_phy_s2_q08_4} Uniform disk of initial mass $M$ from which circular portion of radius $R$ is then removed. M.I. of remaining mass about axis which is perpendicular to plane of plate. & (S) & $\dfrac{MR^2}6$\\ \hline
 & & (T) & $\dfrac{MR^2}8$\\ \hline
\end{tabular}}\par
\begin{qoptions}
\opt{A}{$\mathrm I\to P;\mathrm{II}\to Q;\mathrm{III}\to R;\mathrm{IV}\to T$}
\opt{B}{$\mathrm I\to R;\mathrm{II}\to Q;\mathrm{III}\to S;\mathrm{IV}\to P$}
\opt{C}{$\mathrm I\to Q;\mathrm{II}\to P;\mathrm{III}\to T;\mathrm{IV}\to R$}
\opt{D}{$\mathrm I\to S;\mathrm{II}\to P;\mathrm{III}\to T;\mathrm{IV}\to Q$}
\end{qoptions}
\qn{9} A block suspended from a spring is released from rest when spring is unstratched. `$x$' represents stretch in spring. Select the appropriate graph taking quantities in List-I as $y$-axis.
\qfig{p1_phy_s2_q09_1}
{\small\setlength{\tabcolsep}{3pt}\begin{tabular}{|c|p{0.5\linewidth}|c|c|}
\hline
 & \textbf{List--I} & & \textbf{List--II}\\ \hline
(I) & The KE of block & (P) & 1\\ \hline
(II) & The work done on the block by gravity & (Q) & 2\\ \hline
(III) & The magnitude of work done on the block by spring & (R) & 3\\ \hline
(IV) & The total mechanical energy of block-earth-spring system & (S) & 4\\ \hline
 & & (T) & 5\\ \hline
\end{tabular}}\par
\begin{qoptions}
\opt{A}{$\mathrm I\to R;\mathrm{II}\to P;\mathrm{III}\to Q;\mathrm{IV}\to S$}
\opt{B}{$\mathrm I\to Q;\mathrm{II}\to T;\mathrm{III}\to S;\mathrm{IV}\to R$}
\opt{C}{$\mathrm I\to P;\mathrm{II}\to Q;\mathrm{III}\to R;\mathrm{IV}\to S$}
\opt{D}{$\mathrm I\to T;\mathrm{II}\to Q;\mathrm{III}\to P;\mathrm{IV}\to R$}
\end{qoptions}
\qn{10} In the list-I shown are system of blocks and in list-II are shown ratio of magnitudes of acceleration $\left(\dfrac{a_A}{a_B}\right)$. Match the list-I with list-II. Neglect friction every where ($m_A=5$kg, $m_B=10$kg)\par
{\small\setlength{\tabcolsep}{3pt}\begin{tabular}{|c|c|c|c|}
\hline
 & \textbf{List--I} & & \textbf{List--II}\\ \hline
(I) & \qfigi{p1_phy_s2_q10_1} & (P) & $\dfrac12$\\ \hline
(II) & \qfigi{p1_phy_s2_q10_2} & (Q) & 1\\ \hline
(III) & \qfigi{p1_phy_s2_q10_3} & (R) & $\dfrac34$\\ \hline
(IV) & \qfigi{p1_phy_s2_q10_4} & (S) & $\dfrac53$\\ \hline
 & & (T) & $\dfrac43$\\ \hline
\end{tabular}}\par
\begin{qoptions}
\opt{A}{$\mathrm I\to Q;\mathrm{II}\to R;\mathrm{III}\to T;\mathrm{IV}\to P$}
\opt{B}{$\mathrm I\to Q;\mathrm{II}\to T;\mathrm{III}\to P;\mathrm{IV}\to S$}
\opt{C}{$\mathrm I\to P;\mathrm{II}\to R;\mathrm{III}\to Q;\mathrm{IV}\to S$}
\opt{D}{$\mathrm I\to S;\mathrm{II}\to R;\mathrm{III}\to T;\mathrm{IV}\to Q$}
\end{qoptions}
\sect{SECTION-II}
\qn{1} A block of ice sliding down from the top of the inclined roof of a house (angle of inclination of roof $=30^\circ$ with the horizontal). The highest and lowest points of roof are at heights of 8.1 m and 5.6 m respectively from the ground. At what horizontal distance (in m) from the starting point will the block hit the ground ? Neglect friction.
\qn{2} One end of a rope is fixed to a vertical wall and the other end is pulled by a horizontal force of 20 N. The shape of the flexible rope is shown in the figure. Find its mass (in kg).
\qfig{p1_phy_s3_q02_1}
\qn{3} Three blocks are connected as shown in the figure. The co-efficient of friction at all surfaces is 0.25. If the minimum force required to move the body with constant velocity is $k\,mg$, then find the value of $k$. [Assume that rod does not exert force in vertical direction]
\qfig{p1_phy_s3_q03_1}
\qn{4} A particle of mass `$m$' moves along the quarter section of the circular path whose centre is at the origin. The radius of the circular path is `$a$'. A force $\vec F=y\hat i-x\hat j$ newton acts on the particle, where $x$, $y$ denote the coordinates of position of the particle. If the work done by this force in taking the particle from point A $(a,0)$ to point B$(0,a)$ along the circular path is $-ka^2$ Joules then find the value of $k$.
\qfig{p1_phy_s3_q04_1}
\qn{5} As shown in the figure a person is pulling a mass 1 kg from ground on a fixed rough hemispherical surface upto the top of the hemisphere with the help of a light inextensible string. Find the work done (in J) by tension in the string if radius of hemisphere is 2m and friction coefficient is 0.22. Assume that the block is pulled with negligible velocity.
\qfig{p1_phy_s3_q05_1}
\qn{6} Two particles $P_1$ and $P_2$ of equal masses situated at $(0,0)$ and $(10,0)$ respectively at $t=0$ moving with constant velocities collided head on at point $(4,0)$ after time $t_0$. If the coefficient of restitution is 1, then what is the $x$-coordinate of the centre of mass of the two particles at $t=2t_0$.
\qn{7} Two balls having masses $m$ and $2m$ are fastened to two light strings of same length $\ell$. The other ends of the strings are fixed at O. The strings are kept is same horizontal line and the system is released from rest. The collision between the balls is elastic. If the difference between their maximum height after collision if $(2a+1)\ell/9$ then find $a$.
\qn{8} A cylindrical solid of mass $10^{-2}$ kg and cross-sectional area $10^{-4}$ m$^2$ is moving parallel to its axis (the $X$-axis) with a uniform speed of $10^3$ m/s in the positive direction. At $t=0$, its front face passes the plane $x=0$. The region to the right of this plane is filled with stationary dust particles of uniform density $10^{-3}$ kg/m$^3$. When a dust particle collides with the face of the cylinder, it sticks to its surface. Assuming that the dimensions of the cylinder remain pratically unchanged, and that the dust sticks only to the front face of the cylinder, find the $x$-coordinate (in m) of the cylinder at $t=150$ s. If your answer is in the form of $\alpha^\beta$ meter then find the value of $\beta$.
\partsec{2}{CHEMISTRY}
\sect{SECTION-I(i)}
\qn{1} 23 gm of a gas having density 2.03 gm/litre at at STP when undergoes complete combustion gives 44 gm $\mathrm{CO_2}$ and 27g $\mathrm{H_2O}$. What is/are correct statements regarding the gas -
\begin{qoptions}
\opt{A}{The gas may be $\mathrm{C_2H_5OH}$}
\opt{B}{The gas may be $\mathrm{CH_3{-}O{-}CH_3}$}
\opt{C}{The gas may be $\mathrm{C_2H_6}$}
\opt{D}{The gas contains only carbon and hydrogen}
\end{qoptions}
\qn{2} $2^\circ$ Alkyls are present in
\begin{qoptions}
\opt{A}{\qfigi{p1_chem_s1_q02_A}}
\opt{B}{\qfigi{p1_chem_s1_q02_B}}
\opt{C}{\qfigi{p1_chem_s1_q02_C}}
\opt{D}{$\mathrm{CH_3{-}CH_2{-}Br}$}
\end{qoptions}
\qn{3} If mass of solid residue obtained is 25 gm in each of the following decomposition reaction per 100gm reactant, which of the following reaction has yield greater than 50\% ?\par
(At. wt. : Na = 23, N = 14, Ca = 40, Pb = 208)
\begin{qoptions}
\opt{A}{$\mathrm{2NaHCO_3\to Na_2CO_3\downarrow+CO_2\uparrow+H_2O\uparrow}$}
\opt{B}{$\mathrm{NaNO_3\to NaNO_2\downarrow+\dfrac12O_2\uparrow}$}
\opt{C}{$\mathrm{Ca(NO_3)_2\to CaO\downarrow+2NO_2\uparrow+\dfrac12O_2\uparrow}$}
\opt{D}{$\mathrm{PbCO_3\to PbO\downarrow+CO_2\uparrow}$}
\end{qoptions}
\qn{4} Which of the following statements is/are CORRECT.
\begin{qoptions}
\opt{A}{$\mathrm{SO_4^{2-}=NO_3^-=I_3^-=XeF_5^-}$ (dipole moment)}
\opt{B}{p$\pi$-d$\pi$ bond is present in both $\mathrm{POCl_3}$ and $\mathrm{XeO_3}$}
\opt{C}{MnO is a basic oxide}
\opt{D}{$\mathrm{HClO_4<HClO_3<HClO_2}$ (acidic strength)}
\end{qoptions}
\qn{5} Select the correct option(s) regarding $\mathrm{XeO_3F_2}$
\begin{qoptions}
\opt{A}{Max number of atoms in one plane = 4}
\opt{B}{Number of planes containing max atoms = 2}
\opt{C}{Three $2p_\pi$-$3d_\pi$ bonds}
\opt{D}{O=Xe=O bond angle is lesser than O=Xe=O bond angle in $\mathrm{XeO_2F_4}$}
\end{qoptions}
\qn{6} Correct order of Boling point is/are
\begin{qoptions}
\opt{A}{$\mathrm{H_2<He}$}
\opt{B}{$\mathrm{H_2<D_2}$}
\opt{C}{$\mathrm{H_2O<D_2O}$}
\opt{D}{$\mathrm{NH_3<SbH_3}$}
\end{qoptions}
\sect{SECTION-I(ii)}
\qn{7} Answer the following by appropriately matching the lists based on the information given in the paragraph.\par
{\small\setlength{\tabcolsep}{2pt}\begin{tabular}{|c|p{0.3\linewidth}|c|p{0.38\linewidth}|}
\hline
\multicolumn{2}{|c|}{\textbf{List--I}} & \multicolumn{2}{c|}{\textbf{List--II}}\\ \hline
(P) & 0.5 mole of $\mathrm{SO_2(g)}$ & (1) & Occupy 11.2 L at 1 atm and 273 K\\ \hline
(Q) & 1 gm molecule of $\mathrm{O_2(g)}$ & (2) & Total no. of atoms = $1.5\times N_A$\\ \hline
(R) & 44 gm of $\mathrm{CO_2}$ gas & (3) & $22N_A$ electrons\\ \hline
(S) & 0.5 mole of $\mathrm{O_3(g)}$ & (4) & weighs 32 gm\\ \hline
 & & (5) & Total no. of protons = $16N_A$\\ \hline
\end{tabular}}\par
The CORRECT option is :
\begin{qoptions}
\opt{A}{$P\to1;Q\to2,3,4;R\to4;S\to1$}
\opt{B}{$P\to1,2,4;Q\to4,5;R\to3;S\to1,2$}
\opt{C}{$P\to2;Q\to1,2,5;R\to4;S\to1$}
\opt{D}{$P\to3;Q\to1,2,5;R\to5;S\to1,5$}
\end{qoptions}
\qn{8}
{\small\setlength{\tabcolsep}{2pt}\begin{tabular}{|c|c|c|p{0.5\linewidth}|}
\hline
\multicolumn{2}{|c|}{\textbf{List-- I}} & \multicolumn{2}{c|}{\textbf{List-- II}}\\ \hline
(P) & $\mathrm{PCl_5(g)}$ & (1) & polar specie\\ \hline
(Q) & $\mathrm{C_6H_6}$ & (2) & involves d-orbital in hybridisation\\ \hline
(R) & $\mathrm{SO_3}$ & (3) & specie has atleast 4 atoms in one plane\\ \hline
(S) & $\mathrm{SO_2Cl_2}$ & (4) & specie has all equal bond length\\ \hline
 & & (5) & non-planar specie\\ \hline
\end{tabular}}\par
The correct option is
\begin{qoptions}
\opt{A}{$P\to2,3,5;Q\to1;R\to3,4;S\to1,5$}
\opt{B}{$P\to1,2,3;Q\to1;R\to3,4;S\to1,5$}
\opt{C}{$P\to2,3,5;Q\to3;R\to3,4;S\to1,5$}
\opt{D}{$P\to1,2,3;Q\to1,3;R\to1,2;S\to1,5$}
\end{qoptions}
\qn{9}
{\small\setlength{\tabcolsep}{2pt}\begin{tabular}{|c|c|c|p{0.38\linewidth}|}
\hline
\multicolumn{2}{|c|}{\textbf{List- I} (Compounds)} & \multicolumn{2}{c|}{\textbf{List- II} (Properties)}\\ \hline
(P) & $\mathrm{N_2O_5(s)}$ & (1) & Cation is polar\\ \hline
(Q) & $\mathrm{Cl_2O_6(s)}$ & (2) & Anion is non polar\\ \hline
(R) & $\mathrm{PCl_5(s)}$ & (3) & p$\pi$-p$\pi$ linkage in cation\\ \hline
(S) & $\mathrm{I_2Cl_6(\ell)}$ & (4) & p$\pi$-d$\pi$ linkage in anion\\ \hline
 & & (5) & Both cation and anion are planar\\ \hline
\end{tabular}}\par
Select the incorrect option :-
\begin{qoptions}
\opt{A}{$P\to2;Q\to2;R\to2;S\to5$}
\opt{B}{$P\to3;Q\to4;R\to2;S\to1$}
\opt{C}{$P\to1;Q\to3;R\to2;S\to5$}
\opt{D}{$P\to5;Q\to1;R\to2;S\to2$}
\end{qoptions}
\qn{10}
{\footnotesize\setlength{\tabcolsep}{2pt}\begin{tabular}{|c|p{0.3\linewidth}|c|p{0.28\linewidth}|}
\hline
\multicolumn{2}{|c|}{\textbf{List-I} (Amount of substance)} & \multicolumn{2}{p{0.42\linewidth}|}{\centering\textbf{List-II} (No. of moles of particular atoms in the given substance)}\\ \hline
(P) & $6.022\times10^{24}$ molecules of $\mathrm{Al_2(SO_4)_3{\cdot}3H_2O}$ & (1) & 15 mole O-atoms\\ \hline
(Q) & 90 gm $\mathrm{C_6H_{12}O_6}$ & (2) & 60 mole H-atoms\\ \hline
(R) & 112 litre $\mathrm{SO_3(g)}$ at 1 atm and $0^\circ$C. & (3) & 3 mole O-atoms\\ \hline
(S) & 3.75 mole $\mathrm{N_2O_4(g)}$ & (4) & 30 mole S-atoms\\ \hline
 & & (5) & 150 mole O-atoms\\ \hline
\end{tabular}}\par
The correct option is
\begin{qoptions}
\opt{A}{$P\to2,4,5;Q\to3;R\to1;S\to1$}
\opt{B}{$P\to1;Q\to2,3;R\to1,4;S\to5$}
\opt{C}{$P\to4;Q\to3,4;R\to1;S\to1,2$}
\opt{D}{$P\to2;Q\to4,5;R\to2,4;S\to3$}
\end{qoptions}
\sect{SECTION-II}
\qn{1} Pyrolusite is an ore of manganese and it contains 58\% by mass $\mathrm{MnO_2}$, 15\% inert impurities and the rest is water. When $\mathrm{MnO_2}$ is heated it converts to $\mathrm{Mn_3O_4}$ and oxygen gas is released. A sample of pyrolusite was heated till constant weight is obtained. Calculate the percentage (by wt) of manganese in the residue. (At. wt. of Mn = 55)
\qn{2} Find the sum of molarity of all the ions present in an aqueous solution of 5M $\mathrm{NaNO_3}$ and 3m $\mathrm{BaCl_2}$. The specific gravity of the given solution is 1.665. Assume 100 \% dissociation of each salt.
\qn{3} 2 m solution is formed when 10 gm solute is present in 100 ml solvent. Calculate the molality of solution which is formed due to dissolution of 25 gm solute in 500 ml solvent.
\qn{4} Calculate the density (in gm/ml) of aqueous NaOH solution of which molarity \& (w/w\%) are equal.
\qn{5} DBE of \qfigi{p1_chem_s3_q05_1}
\qn{6} Among the following, total number of planar species is ?\par
(i) $\mathrm{SF_4}$ \quad (ii) $\mathrm{BrF_3}$ \quad (iii) $\mathrm{XeF_2}$ \quad (iv) $\mathrm{IF_5}$ (v) $\mathrm{SbF_4^-}$ \quad (vi) $\mathrm{SF_5^-}$ \quad (vii) $\mathrm{SeF_3^+}$ \quad (viii) $\mathrm{CH_3^+}$ (ix) $\mathrm{PCl_4^+}$
\qn{7} If $x$ = Number of bond angles in $\mathrm{PCl_5}$ which are less than angle in regular tetrahedral .\par
$y$ = Number of equal shorter bond lengths in dichromate ion.\par
$z$ = Total number of bond angles in $\mathrm{P_4}$\par
Then the value of expression $|x+y-z|$ is :-
\qn{8} Which is/are "Pyro" acid :\par
$\mathrm{H_2N_2O_2}$, $\mathrm{H_2S_2O_5}$, $\mathrm{H_2S_2O_7}$, $\mathrm{HClO_4}$, $\mathrm{H_4P_2O_5}$, $\mathrm{H_4P_2O_7}$, $\mathrm{H_2SO_5}$
\partsec{3}{MATHEMATICS}
\sect{SECTION-I(i)}
\qn{1} The equation of the straight line equally inclined to the coordinate axes and equidistant from the points $(2,0)$ and $(0,4)$ can be-
\begin{qoptions}
\opt{A}{$y=x-1$}
\opt{B}{$y=x+1$}
\opt{C}{$x+y=3$}
\opt{D}{$x+y+3=0$}
\end{qoptions}
\qn{2} Let $A(0,3)$, $B(-2,0)$ and $C(6,1)$ be the vertices of a triangle. If a moving point $M(0,\beta+1)$ lies within the $\Delta ABC$, then possible value(s) of $\beta$ is/are-
\begin{qoptions}
\opt{A}{$\dfrac75$}
\opt{B}{$-\dfrac34$}
\opt{C}{$-\dfrac{11}{20}$}
\opt{D}{$\dfrac12$}
\end{qoptions}
\qn{3} If sides $AB$ and $AC$ of a $\Delta ABC$ are given by equations $3x+y=4$ and $x+3y=4$ respectively and mid point of $BC$ is $(2,2)$, then which of the following is/are true :-
\begin{qoptions}
\opt{A}{Equation of $BC$ is $x-y=4$}
\opt{B}{Equation of $BC$ is $x+y=4$}
\opt{C}{Area of $\Delta ABC$ is $4$}
\opt{D}{Area of $\Delta ABC$ is $8$}
\end{qoptions}
\qn{4} If $T_n=\begin{vmatrix}n-1&n-2&n-6\\2n-4&2n-6&2n-11\\3n-9&3n-12&3n-18\end{vmatrix}$ where $n\in N$, then which of the following is/are true ?
\begin{qoptions}
\opt{A}{$\displaystyle\prod_{n=1}^{n}T_n=6^n$}
\opt{B}{$\displaystyle\sum_{n=1}^{10}T_n=60$}
\opt{C}{$\dfrac{T_{n+1}}{T_n}=2$}
\opt{D}{$T_{100}\cdot T_{101}=48$}
\end{qoptions}
\qn{5} Let $P$ and $Q$ be any two points on the lines represented by $2x-3y=0$ and $2x+3y=0$ respectively. If area of triangle $OPQ$, ($O$ is origin) is $5$, then equation of locus of mid point of $PQ$ can be
\begin{qoptions}
\opt{A}{$4x^2-9y^2+30=0$}
\opt{B}{$9x^2-4y^2+30=0$}
\opt{C}{$9x^2-4y^2-30=0$}
\opt{D}{$4x^2-9y^2-30=0$}
\end{qoptions}
\qn{6} Let $f(x)=ax^2+bx+c$; $a,b,c\in R$ and equation $f(x)=0$ has exactly one solution in interval $(m,n)$, then $f(m)\cdot f(n)$ can be equal to -
\begin{qoptions}
\opt{A}{$\sin1\cos2$}
\opt{B}{$\sin2\cos1$}
\opt{C}{$\sin2\cos2$}
\opt{D}{$\sin3\cos2$}
\end{qoptions}
\sect{SECTION-I(ii)}
\qn{7} Consider the quadratic equation $(a^2-3a+2)x^2+(2a^3-a^2-11a+10)x+(2a^2+3a-5)=0$ then find the sum of all the value(s) of $a$ (where $a\in R$) for which the given equation has\par
{\small\setlength{\tabcolsep}{2pt}\begin{tabular}{|c|p{0.4\linewidth}|c|c|}
\hline
\multicolumn{2}{|c|}{\textbf{List I}} & \multicolumn{2}{c|}{\textbf{List II}}\\ \hline
(I) & both roots infinity & (P) & $-7$\\ \hline
(II) & both roots zero & (Q) & $2$\\ \hline
(III) & exactly one root infinity & (R) & $\dfrac{-5}{2}$\\ \hline
(IV) & both roots are multiplicative inverse of each other & (S) & $\phi$\\ \hline
 & & (T) & $0$\\ \hline
\end{tabular}}\par
\begin{qoptions}
\opt{A}{$I\to R;\allowbreak II\to Q;\allowbreak III\to P;\allowbreak IV\to S$}
\opt{B}{$I\to Q;\allowbreak II\to R;\allowbreak III\to P;\allowbreak IV\to S$}
\opt{C}{$I\to Q;\allowbreak II\to R;\allowbreak III\to S;\allowbreak IV\to P$}
\opt{D}{$I\to R;\allowbreak II\to Q;\allowbreak III\to S;\allowbreak IV\to P$}
\end{qoptions}
\qn{8} Match List-I with List-II and select the correct answer using the code given below the list.\par
{\small\setlength{\tabcolsep}{2pt}\begin{tabular}{|c|p{0.4\linewidth}|c|c|}
\hline
\multicolumn{2}{|c|}{\textbf{List-I}} & \multicolumn{2}{c|}{\textbf{List-II}}\\ \hline
(I) & Distance between the two lines $x+\sqrt3y-2=0$ and $x+\sqrt3y+2=0$ is & (P) & $1$\\ \hline
(II) & The variable line $x+\lambda y+\lambda-5=0$ ($\lambda\in R$) passes through a fixed point. The $x$ co-ordinate of fixed point is & (Q) & $2$\\ \hline
(III) & The $x$-coordinates of the point of intersection of lines $3x+4y=9$ and $y=mx+1$ is an integer, if $m$ is also an integer then sum of possible values of $m$ is & (R) & $5$\\ \hline
(IV) & The $x$-coordinate of incentre of triangle, whose vertices are $(1,\sqrt3)$, $(0,0)$ and $(2,0)$, is & (S) & $-3$\\ \hline
 & & (T) & $-1$\\ \hline
\end{tabular}}\par
\begin{qoptions}
\opt{A}{$I\to Q;\allowbreak II\to R;\allowbreak III\to P;\allowbreak IV\to S$}
\opt{B}{$I\to Q;\allowbreak II\to R;\allowbreak III\to S;\allowbreak IV\to P$}
\opt{C}{$I\to R;\allowbreak II\to Q;\allowbreak III\to S;\allowbreak IV\to P$}
\opt{D}{$I\to P;\allowbreak II\to Q;\allowbreak III\to S;\allowbreak IV\to R$}
\end{qoptions}
\qn{9} Match List-I with List-II and select the correct answer using the code given below the list.\par
{\small\setlength{\tabcolsep}{2pt}\begin{tabular}{|c|p{0.4\linewidth}|c|c|}
\hline
\multicolumn{2}{|c|}{\textbf{List-I}} & \multicolumn{2}{c|}{\textbf{List-II}}\\ \hline
(I) & Number of common integers in two arithmetic progression $1,3,5,7,\ldots157$ and $2,5,8,11\ldots161$ is & (P) & $0$\\ \hline
(II) & If $a$ and $b$ are roots of $11x^2-4x-2=0$, then value of $10(1+a+a^2+\ldots\infty)(1+b+b^2+\ldots\infty)$ is & (Q) & $26$\\ \hline
(III) & If the sum of first $20$ terms is same as the sum of first $10$ terms in an arithmetic progression then the sum of first $30$ terms is & (R) & $22$\\ \hline
(IV) & If sum of {\footnotesize$\dfrac{3}{1^2\cdot2^2}+\dfrac{5}{2^2\cdot3^2}+\dfrac{7}{3^2\cdot4^2}+\ldots+\dfrac{199}{99^2\cdot100^2}$} is $\dfrac ab$ (where $a$ and $b$ are coprime) the value of $20|b-a|$ is & (S) & $20$\\ \hline
 & & (T) & $10$\\ \hline
\end{tabular}}\par
\begin{qoptions}
\opt{A}{$I\to Q;\allowbreak II\to P;\allowbreak III\to P;\allowbreak IV\to R$}
\opt{B}{$I\to Q;\allowbreak II\to R;\allowbreak III\to P;\allowbreak IV\to S$}
\opt{C}{$I\to P;\allowbreak II\to Q;\allowbreak III\to P;\allowbreak IV\to Q$}
\opt{D}{$I\to P;\allowbreak II\to S;\allowbreak III\to R;\allowbreak IV\to P$}
\end{qoptions}
\qn{10} A triangle $ABC$ is shown in the figure with $\angle C=60^\circ$, $a=5$ \& $b=4$.\par
\qfig{p1_math_s2_q10_1}
Match List-I with List-II and select the correct answer using the code given below the list.\par
{\small\setlength{\tabcolsep}{2pt}\begin{tabular}{|c|p{0.4\linewidth}|c|c|}
\hline
\multicolumn{2}{|c|}{\textbf{List-I}} & \multicolumn{2}{c|}{\textbf{List-II}}\\ \hline
(I) & The value of the side opposite to angle $\angle C$ is & (P) & $5\sqrt3$\\ \hline
(II) & If $R$ is the circumradius of the triangle $ABC$, the value of $R^2$ is & (Q) & $7$\\ \hline
(III) & Length of internal angle bisector of $C$ is & (R) & $\sqrt{21}$\\ \hline
(IV) & Area of triangle $ABC$ is & (S) & $\dfrac{20}{3\sqrt3}$\\ \hline
 & & (T) & $3\sqrt5$\\ \hline
\end{tabular}}\par
\begin{qoptions}
\opt{A}{$I\to R;\allowbreak II\to Q;\allowbreak III\to S;\allowbreak IV\to P$}
\opt{B}{$I\to Q;\allowbreak II\to R;\allowbreak III\to S;\allowbreak IV\to P$}
\opt{C}{$I\to P;\allowbreak II\to Q;\allowbreak III\to R;\allowbreak IV\to S$}
\opt{D}{$I\to R;\allowbreak II\to S;\allowbreak III\to P;\allowbreak IV\to Q$}
\end{qoptions}
\sect{SECTION-II}
\qn{1} If $3\sqrt3$, $2\sqrt2\sin^2\theta$, $\sin\theta$, ($\theta\neq n\pi$, $n\in I$) are in geometric progression, then number of values of $\theta$ in $[-2\pi,3\pi]$ is
\qn{2} A rod of unit length slides along co-ordinate axes such that its two ends $A$ and $B$ always lie on positive $X$ and $Y$ axes respectively. Then, the locus of centroid of triangle $OAB$ (where $O$ is origin), is given by $k\left(x^2+y^2\right)=1$; then the value of $k$ is
\qn{3} If number of possible triplet $(x,y,z)$, which satisfies the given determinant, where $x,y,z\in I^+$\par
$\begin{vmatrix}x^3+1&x^2y&x^2z\\xy^2&y^3+1&y^2z\\xz^2&yz^2&z^3+1\end{vmatrix}=11$ is $k$ then value of $\dfrac k2$ is
\qn{4} Let $x$ and $y$ be positive real numbers and $\theta$ an angle such that $\theta\neq\dfrac{n\pi}2$ for any integer $n$. Suppose $\dfrac{\sin\theta}{x}=\dfrac{\cos\theta}{y}$ and $\dfrac{\cos^4\theta}{x^4}+\dfrac{\sin^4\theta}{y^4}=\dfrac{97\sin2\theta}{x^3y+xy^3}$, then the value of $\dfrac xy+\dfrac yx$ is
\qn{5} If $2,h_1,h_2,h_3\ldots h_{20},5$ are in H.P., $2,a_1,a_2,a_3\ldots a_{20},5$ are in A.P. and $2,g_1,g_2,g_3\ldots g_{20},5$ are in G.P. and $L=\displaystyle\prod_{i=1}^{20}a_ih_ig_i$, then $\left(\dfrac{\log_{10}L}{10}\right)$ is equal to
\qn{6} If a point $P$ moves such that its distance from line $y=\sqrt3x-7$ is same as its distance from $\left(3\sqrt3,\ 2\right)$, then the area of the curve, described by $p$, enclosed between the coordinate axes is $\dfrac{\sqrt3}{2}k$ (in square units), then the value of $k$ is
\qn{7} Let $a_2,a_3\in R$ such that $|a_2-a_3|=6$ and $f(x)=\begin{vmatrix}1&a_3&a_2\\1&a_3&2a_2-x\\1&2a_3-x&a_2\end{vmatrix}$, $x\in R$. Then the greatest value of $f(x)$ is
\qn{8} Let the expression $E=8^a+8^b-3.2^{a+b}$ takes its minimum value '$p$' at $a=\alpha$ \& $b=\beta$, then perpendicular distance of the point $P(\alpha,\beta)$ from the line $x+y+2p=0$ is $\sqrt\lambda$, then the value of $\lambda$ is
\papersec{2}
\partsec{1}{PHYSICS}
\sect{SECTION-I(i)}
\qn{1} A uniform disc of mass $2$ kg and radius $1$m is mounted on a horizontal axle supported on fixed frictionless bearings. A light chord is wrapped around the rim of the disc and mass of $1$ kg is tied to the free end. If it is released from rest, which of the following is INCORRECT ? (take $g=10\ \mathrm{m/s^2}$)
\begin{qoptions}
\opt{A}{The tension in the chord is $5$ N}
\opt{B}{In first four seconds the angular displacement of the disc is $50$ rad}
\opt{C}{The work done by the torque on the disc in first four seconds is $200$ J}
\opt{D}{The increase in the kinetic energy of the mass in the first four seconds in $200$ J}
\end{qoptions}
\qn{2} A light flexible thread passes over a small frictionless pulley. Two blocks of masses $1$ kg and $3$ kg are attached with a thread and heavy block rests on a surface. A particle of mass $1$ kg moving upwards with a velocity $10$ m/s collides with the lighter block elastically at $t=0$. ($g=10\ \mathrm{m/s^2}$). Choose the INCORRECT option.\par
\qfig{p2_phy_s1_q02_1}
\begin{qoptions}
\opt{A}{Velocity with which heavy block is jerked into motion is $2.5$ m/s}
\opt{B}{Maximum height ascended by heavy block after it is jerked into motion is $0.625$ m}
\opt{C}{impulse of tension in jerk is $7.5$ Ns}
\opt{D}{Time after which block again strikes the surface is $2.5$ s.}
\end{qoptions}
\qn{3} In the figure, the block $B$ of mass $m$ starts from at rest top of a wedge $W$ of mass $M$. All surfaces are without friction. $W$ can slide on the ground. $B$ slides down onto the ground, moves along it with a speed $v$, has an elastic collision with the wall, and climbs back onto $W$.\par
\qfig{p2_phy_s1_q03_1}
\begin{qoptions}
\opt{A}{$B$ will reach the top of $W$ again.}
\opt{B}{From the beginning, till the collision with the wall, the centre of mass of '$B$' plus '$W$' is stationary.}
\opt{C}{After the collision, centre of mass of '$B$ plus $W$' moves with the horizontal component of velocity $\dfrac{2mv}{m+M}$.}
\opt{D}{When $B$ reaches its highest position on $W$, the speed of $W$ is $\dfrac{mv}{m+M}$.}
\end{qoptions}
\qn{4} A small bead of mass $m=4.0$ kg is threaded on a homogeneous frictionless ring of mass $M=6.0$ kg and radius $R=0.5$ m. Initially the system rests in free space. Now the bead is struck to impart a velocity $u=5.0$ m/s along the tangent to the ring. Find minimum kinetic energy of the bead in subsequent motion.
\begin{qoptions}
\opt{A}{$1$J}
\opt{B}{$2$J}
\opt{C}{$\dfrac{16}{3}$J}
\opt{D}{$8$J}
\end{qoptions}
\sect{SECTION-I(ii)}
\qn{5} A semicircular rigid cylindrical tube is placed on a horizontal floor. Two identical balls that loosely fit inside the tube are simultaneously released from the top of the tube. Each ball moves down the tube on either side. Which of the following statements best describe the total force of normal reaction between the tube and the floor ?\par
\qfig{p2_phy_s2_q05_1}
\begin{qoptions}
\opt{A}{It decreases continually.}
\opt{B}{It increases continually.}
\opt{C}{It first decreases, acquires a minimum value then increases.}
\opt{D}{The tube may leave the ground.}
\end{qoptions}
\qn{6} A light inextensible cord is wrapped around a frictionless cylinder of radius $R$ fixed vertically on a floor that is lubricated with oil. A bead of mass $m$ is tied at one end of the cord and the other end $P$ is free as shown by the top view of the situation. Force of viscous drag on the bead due to layer of the lubricating oil is proportional to the speed of the bead and the proportionality constant is $k$. If the free end of the cord is pulled with a constant speed $u$ along the cylinder, the bead eventually acquires a steady motion on a circular path for speed of the bead ($v$) and radius of its circular path $r$ in its steady motion are :\par
\qfig{p2_phy_s2_q06_1}
\begin{qoptions}
\opt{A}{$v=u\sqrt{1+\left(\dfrac{mu}{kR}\right)^2}$}
\opt{B}{$v=\dfrac{mu^2}{kR}$}
\opt{C}{$r=\dfrac{mu}{k}$}
\opt{D}{$r=R\sqrt{1+\left(\dfrac{mu}{kR}\right)^2}$}
\end{qoptions}
\qn{7} A triangular block is moving in its own plane such that at any instant speeds of corner $A$ and $B$ are $v$ each, as shown in figure. Then\par
\qfig{p2_phy_s2_q07_1}
\begin{qoptions}
\opt{A}{Velocity of end $C$ is zero.}
\opt{B}{Velocity of end $C$ is $v$ upwards.}
\opt{C}{Velocity of mid-point of side $BC$ is zero.}
\opt{D}{Speed of mid-point of side $AC$ is $v/2$.}
\end{qoptions}
\qn{8} In a stationary cart of mass $32m$ and length $d$, a small block of mass $m$ is projected along its length with velocity $v$ towards front end from rear end as shown in figure. If $e=\dfrac12$ and friction is absent then :-\par
\qfig{p2_phy_s2_q08_1}
\begin{qoptions}
\opt{A}{final velocity of cart is $\dfrac{v}{33}$}
\opt{B}{after five collisions block will be at rest w.r.t. ground.}
\opt{C}{time at which the block comes to rest w.r.t. ground is $\dfrac{31d}{v}$}
\opt{D}{time interval between fourth and fifth collision is $\dfrac{16d}{v}$.}
\end{qoptions}
\qn{9} Three identical discs $A$, $B$ and $C$ rest on a smooth horizontal plane as shown in figure. The disc $A$ is set in motion with velocity $v$ along the perpendicular bisector of line $BC$. The distance between disc $B$ and $C$ is $n$ times the diameter of each disc. For which of the following values of $n$ will the disc $A$ continue moving after elastic collision ?\par
\qfig{p2_phy_s2_q09_1}
\begin{qoptions}
\opt{A}{$n=\sqrt{2.5}$}
\opt{B}{$n=\sqrt2$}
\opt{C}{$n=\sqrt{1.5}$}
\opt{D}{$n=\sqrt3$}
\end{qoptions}
\qn{10} A particle of unit mass is moving along $x$-axis under the influence of a force and its total energy is conserved. If its potential energy changes as\par
$U_4(x)=\dfrac{U_0}{2}\left[\dfrac xa-\dfrac13\left(\dfrac xa\right)^3\right]$\par
Which of the following are correct ?
\begin{qoptions}
\opt{A}{The force acting on the particle is zero at $x=a$.}
\opt{B}{The force acting on the particle is zero at $x=0$}
\opt{C}{The particle experiences an attractive force towards $x=0$ in the region $|x|<a$.}
\opt{D}{The force acting on the particle is zero at $x=-a$}
\end{qoptions}
\sect{SECTION-III}
\qn{1} A safety device brings the blade of a power mower from an initial angular speed of $\omega_1$ to rest in $1.00$ revolution. At the same constant acceleration, how many revolutions would it take the blade to come to rest from an initial angular speed $\omega_3$ that was three times as great, $\omega_3=3\omega_1$ ?
\qn{2} In the arrangement shown in the figure, $m_1=3$kg, $m_2=2$kg, $M=5$ kg and there is no friction at any contact surface. The acceleration due to gravity is $10\ \mathrm{m/s^2}$. The system is released from rest. After $1$ second, the velocity of $m_1$ is $3$ m/s. Kinetic energy of the system at this instant is $k$, then $\dfrac{k}{27}$ will be round off to nearest integer.\par
\qfig{p2_phy_s3_q02_1}
\qn{3} A typical flying insect applies an average force equal to twice its weight during each downward stroke while hovering. Take the mass of the insect to be $10$ g and assume the wings move an average downward distance of $1.0$ cm during each stroke. Assuming $100$ downward strokes per second. If the average power output of the insect is $P$, then $\dfrac1P$ will be :
\qn{4} A uniform, solid disk with mass $m$ and radius $R$ is pivoted about a horizontal axis through its center. A small object of the same mass $m$ is glued to the rim of the disk. If the disk is released from rest with the small object at the end of a horizontal radius, find the angular speed when the small object is directly below the axis. (Take $R=\dfrac{10}{3}$m)
\qn{5} From the top of a tower of height $65$ m, a ball is projected with velocity $40$ m/s at an angle of elevation of $53^\circ$. The time (in seconds) after which its velocity becomes perpendicular to its initial velocity (take $g=10\ \mathrm{m/s^2}$) is $t$ and its height from the ground at that instant is $h$, find the value of $\dfrac{ht}{125}$.
\qn{6} Velocity versus displacement of a particle moving in a straight line is $v=\sqrt{9+4s}$, here $v$ is in m/s and $s$ in meter. Mass of the particle is $2$kg. The time in seconds at which average power is $\dfrac34$th of the instantaneous power is $\dfrac n2$. $n$ is
\qn{7} A block of mass $4$kg is moving with a velocity of $7$ m/sec on a surface. It collides with another block of mass $3$kg elastically. The surface is smooth for $4$ kg block but rough for $3$kg block ($\mu=0.2$). Find the time (in sec) after which next collision will occur ($g=10\ \mathrm{m/s^2}$)
\qn{8} A particle of mass $1$ kg moves with a constant speed $v=1$ m/s along the curve $ACB$ in vertical plane whose equation is $y=0.5\cos\dfrac{2\pi x}{\ell}$, where $x$ and $y$ are in metre and $\ell=1$ m. The normal reaction on particle at lowest point $B$ is $(5n)$ N. (take $\pi^2=10$ and $g=10\ \mathrm{m/s^2}$). The value of $n$ is\par
\qfig{p2_phy_s3_q08_1}
\partsec{2}{CHEMISTRY}
\sect{SECTION-I(i)}
\qn{1} How many primary carbon and secondary hydrogens are present\par
\qfig{p2_chem_s1_q01_1}
\begin{qoptions}
\opt{A}{$15,0$}
\opt{B}{$5,0$}
\opt{C}{$0,15$}
\opt{D}{$0,5$}
\end{qoptions}
\qn{2} Which of the following compound has all $1^\circ$, $2^\circ$, $3^\circ$ \& $4^\circ$ carbons.
\begin{qoptions}
\opt{A}{Iso octane}
\opt{B}{Iso nonane}
\opt{C}{Iso hexane}
\opt{D}{Neo hexane}
\end{qoptions}
\qn{3} Select the CORRECT order of bond angle.
\begin{qoptions}
\opt{A}{$\mathrm{NO_2^+>NO_2^->NO_3^-}$}
\opt{B}{$\mathrm{OCl_2>OF_2>H_2O}$}
\opt{C}{$\mathrm{BF_3=BCl_3>CCl_4}$}
\opt{D}{$\mathrm{NCl_3<PCl_3<AsCl_3}$}
\end{qoptions}
\qn{4} Which of the following order is INCORRECT for boiling point of given molecules --
\begin{qoptions}
\opt{A}{n-pentane $>$ neo-pentane}
\opt{B}{$\mathrm{CH_4<NH_3<H_2O<HF}$}
\opt{C}{$\mathrm{C_2H_6<C_2F_6}$}
\opt{D}{$\mathrm{H_3BO_3>(CH_3)_3BO_3}$}
\end{qoptions}
\sect{SECTION-I(ii)}
\qn{5} A gas mixture contains $1$ mol $\mathrm{CH_4}$ and $2$ mol $\mathrm{SO_3}$. Correct statement(s) is / are -
\begin{qoptions}
\opt{A}{Average molar mass of the mixture is $\dfrac{176}{6}$}
\opt{B}{\% decrease in average molar mass of the mixture if $\mathrm{SO_3}$ is completely dissociated in $\mathrm{SO_2}$ and $\mathrm{O_2}$ = $25\%$}
\opt{C}{Mass ratio of $\mathrm{CH_4}$ and $\mathrm{SO_3}$ initially present in the mixture is $1:10$}
\opt{D}{If $1$ mol more $\mathrm{CH_4}$ is added in the given gaseous mixture then new molar mass of the mixture is $48\ \mathrm{gm\,mol^{-1}}$.}
\end{qoptions}
\qn{6} What is the molality of a solution made by dissolving $156$ g of bromothymol blue ($\mathrm{C_{27}H_{28}Br_2O_5S}$) in $1.00$ L of ethanol on a winter's day at $10^\circ$C ? The density of ethanol at this temperature is $0.80\ \mathrm{kg\,L^{-1}}$ : (Br = 80, S = 32)
\begin{qoptions}
\opt{A}{$0.1000\ \mathrm{mol\,kg^{-1}}$}
\opt{B}{$0.1284\ \mathrm{mol\,kg^{-1}}$}
\opt{C}{$0.1602\ \mathrm{mol\,kg^{-1}}$}
\opt{D}{$0.3125\ \mathrm{mol\,kg^{-1}}$}
\end{qoptions}
\qn{7} Consider the following three orbitals:\par
\qfig{p2_chem_s2_q07_3}\qfig{p2_chem_s2_q07_2}\qfig{p2_chem_s2_q07_1}
Correct statement(s) regarding given information is/are
\begin{qoptions}
\opt{A}{Orbitals (i) and (ii) can never form any type of covalent bond if $x$ is internuclear axis}
\opt{B}{If internuclear axis is '$x$', then combination of (ii) and (iii) orbitals can form $\pi-$bond}
\opt{C}{Orbital (iii) can form $\delta-$bond with other orbital having same orientation of lobes.}
\opt{D}{If internuclear axis is '$x$', then combination of (i) and (iii) orbitals can form $\pi-$bond}
\end{qoptions}
\qn{8} Which of the following orders incorrect between $\mathrm{PCl_3F_2}$ and $\mathrm{PCl_2F_3}$ :-
\begin{qoptions}
\opt{A}{$\mathrm{PCl_3F_2>PCl_2F_3}$ (Dipole moment)}
\opt{B}{$\mathrm{PCl_3F_2>PCl_2F_3}$ (Total number of lone pair)}
\opt{C}{$\mathrm{PCl_3F_2>PCl_2\ F_3}$ (Cl-P-Cl Bond angle)}
\opt{D}{$\mathrm{PCl_3F_2>PCl_2F_3}$ (P--Cl Bond length)}
\end{qoptions}
\qn{9} A solution is made by dissolving $\mathrm{CaBr_2}$ in water (solvent) such that mass fraction of solute and solvent is same in the solution. The molality of solution is -\par
(At. wt. : Br = 80, Ca = 40)
\begin{qoptions}
\opt{A}{$2.5$ m}
\opt{B}{$55.55$ m}
\opt{C}{$2$ m}
\opt{D}{$5$ m}
\end{qoptions}
\qn{10} Given the correct order of initials 'T' or 'F' for the following statements. Use 'T' if statement is true and 'F' if it is false\par
(I) In $\mathrm{NH_3}$, the orbital in which lone pair is present has exactly $25\%$ s-character\par
(II) In $\mathrm{F_2O_2}$, polar as well as non polar both bonds are present\par
(III) In $\mathrm{IO_2F_2^-}$, one of the $\pi$-bond can be formed by $d_{z^2}$ orbital of central atom\par
(IV) In $\mathrm{HCO_3^-}$, all C--O bonds are not same
\begin{qoptions}
\opt{A}{T T T T}
\opt{B}{F T F T}
\opt{C}{F T T T}
\opt{D}{T F T F}
\end{qoptions}
\sect{SECTION-III}
\qn{1} $100$ ml aqueous solution containing equimolar mixture of $\mathrm{Ca(OH)_2}$ and $\mathrm{Al(OH)_3}$ requires $0.5$ litre of $4$M HCl for complete neutralisation. Molarity of $\mathrm{Ca(OH)_2}$ in the original solution is (Assume $\mathrm{Ca(OH)_2}$ and $\mathrm{Al(OH)_3}$ completely soluble in water)
\qn{2} $500$ ml of $2$M $\mathrm{CH_3COOH}$ solution is mixed with $600$ ml $12\%$ w/v $\mathrm{CH_3COOH}$ solution, then calculate the final molarity of solution.
\qn{3} A $1800$ gm mixture of anhydrous $\mathrm{CuSO_4(s)}$ and its hydrated form [$\mathrm{CuSO_4{\cdot}5H_2O(s)}$] undergoes $20\%$ loss in mass in heating. Mole fraction of $\mathrm{CuSO_4(s)}$ in mixture is (Atomic mass of Cu = 64) (Assume the formation of anhydrous salt on heating)\par
If the answer is "X", report answer as 9X
\qn{4} $90$ g of a silver coin was dissolved in strong nitric acid, and excess of sodium chloride solution was added. The silver chloride precipitate was dried and weighed $7.175$ g. Calculate the percentage of silver in the coin (Atomic mass of Ag = 108)\par
$\mathrm{Ag+2HNO_3\longrightarrow AgNO_3+NO_2+H_2O}$\par
$\mathrm{AgNO_3+NaCl\longrightarrow AgCl+NaNO_3}$
\qn{5} Find the number of molecule which has p$\pi$--d$\pi$ bonds\par
$\mathrm{CO_2}$, $\mathrm{C_2H_2}$, $\mathrm{SO_2}$, $\mathrm{NO_2^-}$, $\mathrm{P_4O_{10}}$, $\mathrm{SO_3}$, $\mathrm{N_2O_5}$, $\mathrm{N_2O_3}$, $\mathrm{Cl_2O_7}$
\qn{6} The number of correct orders wrt the property mentioned along with is:\par
(i) $\mathrm{HF>HCl>HBr>HI}$ (Boiling point)\par
(ii) $\mathrm{H_2S<H_2Se<H_2Te}$ (Acidic strength)\par
(iii) $\mathrm{Cl_2>F_2}$ (Bond energy)\par
(iv) $\mathrm{NH_3>PH_3>AsH_3>SbH_3}$ (Proton affinity)\par
(v) $\mathrm{MnO_2>Mn_2O_7}$ (Acidic strength)\par
(vi) $\mathrm{HOCl>HClO_3}$ (Acidic strength)\par
(vii) $\mathrm{H_2O(s)<D_2O(s)}$ (Density)\par
(viii) $\mathrm{CH_3CH_2OH<CH_3CHO<CH_3COOH}$ (Boiling point)\par
(ix) $\mathrm{H_2O_2>H_2O}$ (Strength of H-bonding)
\qn{7} Consider the following species.\par
$\mathrm{SiH_4}$, $\mathrm{BaTiSi_3O_9}$, $\mathrm{NO_2^+}$, $\mathrm{CO_2}$, $\mathrm{C_6H_6}$\par
Find the value of $(x+y)$, if\par
$x$ = number of species containing polar covalent bonds\par
$y$ = number of species containing $\pi$-bonds
\qn{8} How many compounds are 'Homocyclic hydrocarbon' among following compounds :\par
\qfig{p2_chem_s3_q08_1}
\partsec{3}{MATHEMATICS}
\sect{SECTION-I(i)}
\qn{1} If both the roots of the equation $x^2-6ax+2-2a+9a^2=0$ exceed $3$, then :
\begin{qoptions}
\opt{A}{$a<1$}
\opt{B}{$a>11/9$}
\opt{C}{$a>3/2$}
\opt{D}{$a<5/2$}
\end{qoptions}
\qn{2} The system of equation\par
$-2x+y+z=1$,\par
$x-2y+z=-2$,\par
$x+y+\lambda z=4$\par
will have no solution if :
\begin{qoptions}
\opt{A}{$\lambda=-2$}
\opt{B}{$\lambda=-1$}
\opt{C}{$\lambda=3$}
\opt{D}{None}
\end{qoptions}
\qn{3} The opposite vertices of a square are $(1,2)$ and $(3,8)$, then equation of diagonal of the square not passing through $(1,2)$ is
\begin{qoptions}
\opt{A}{$3x-y-1=0$}
\opt{B}{$3y+x=17$}
\opt{C}{$3x+y+1=0$}
\opt{D}{$3y+x=1$}
\end{qoptions}
\qn{4} The line $5x+3y=30$ meet the $x$-axis at $A$ and $y$ axis at $B$. The perpendicular bisector of $AB$ meet the line $x-y=0$ at $P$. The area of $\Delta PAB$ is :
\begin{qoptions}
\opt{A}{$30$ sq. units}
\opt{B}{$32$ sq. units}
\opt{C}{$34$ sq. units}
\opt{D}{$36$ sq. units}
\end{qoptions}
\sect{SECTION-I(ii)}
\qn{5} If $A_i$, $G_i$ and $H_i$ ($i=1,2,\ldots10$) are respectively $10$ AM's, GM's and HM's between two distinct positive numbers, then -
\begin{qoptions}
\opt{A}{$\dfrac{A_1}{A_2}=\dfrac{H_9}{H_{10}}$}
\opt{B}{$(A_1H_{10})^5=G_1\cdot G_2\cdot G_3\ldots G_{10}$}
\opt{C}{$(A_1+A_2+\ldots+A_{10})=\dfrac1{H_1}+\dfrac1{H_2}+\ldots..+\dfrac1{H_{10}}$}
\opt{D}{$A_1A_2=H_9\cdot H_{10}$}
\end{qoptions}
\qn{6} If $\begin{vmatrix}2bc-a^2&c^2&b^2\\c^2&2ca-b^2&a^2\\b^2&a^2&2ab-c^2\end{vmatrix}=(a^3+b^3+c^3+kabc)^2$, then $k$ is equal to
\begin{qoptions}
\opt{A}{$2$}
\opt{B}{$-2$}
\opt{C}{$3$}
\opt{D}{$-3$}
\end{qoptions}
\qn{7} Three lines $px+qy+r=0$, $qx+ry+p=0$ and $rx+py+q=0$ are concurrent if :-
\begin{qoptions}
\opt{A}{$p+q+r=0$}
\opt{B}{$p^2+q^2+r^2=pq+qr+rp$}
\opt{C}{$p^3+p^3+r^3=3\,pqr$}
\opt{D}{None of these}
\end{qoptions}
\qn{8} If $a^2+9b^2-4c^2=6\,ab$ then the family of lines $ax+by+c=0$ are concurrent at :-
\begin{qoptions}
\opt{A}{$(1/2,\,3/2)$}
\opt{B}{$(-1/2,\,-3/2)$}
\opt{C}{$(-1/2,\,3/2)$}
\opt{D}{$(1/2,\,-3/2)$}
\end{qoptions}
\qn{9} $\sin^2x+2\sin x\cos x-3\cos^2x=0$ if
\begin{qoptions}
\opt{A}{$\tan x=3$}
\opt{B}{$\tan x=-1$}
\opt{C}{$x=n\pi+\pi/4$, $n\in I$}
\opt{D}{$x=n\pi+\tan^{-1}(-3)$, $n\in I$}
\end{qoptions}
\qn{10} If $\begin{vmatrix}1&4&20\\1&-2&5\\1&2x&5x^2\end{vmatrix}=0$, then values of $x$ satisfying the equation is/are
\begin{qoptions}
\opt{A}{$2$}
\opt{B}{$-2$}
\opt{C}{$1$}
\opt{D}{$-1$}
\end{qoptions}
\sect{SECTION-III}
\qn{1} The total number of distinct $x\in R$ for which $\begin{vmatrix}x&x^2&1+x^3\\2x&4x^2&1+8x^3\\3x&9x^2&1+27x^3\end{vmatrix}=10$ is
\qn{2} Maximum value of determinant {\footnotesize$\begin{vmatrix}1&\sin\theta&1\\-\sin\theta&1&\sin\theta\\-1&-\sin\theta&1\end{vmatrix}$} is
\qn{3} If $a$, $2b$, $c$ are distinct numbers in A.P. and line $ax+by+c=0$ always passes through a fixed point $(p,q)$, then value of $|p+2q|$ is :-
\qn{4} In $\Delta ABC$, if $A\equiv(1,2)$ and equations of the medians through $B$ and $C$ are $x+y=5$ and $x=4$ respectively, then point $B$ must be $(h,k)$, then $2h+3k$ equals
\qn{5} If $(\alpha,\beta)$ be the coordinates of incentre of the triangle whose vertices are $(0,0)$, $(6,0)$ and $(0,8)$, then the value of $\alpha+\beta$ must be
\qn{6} Given that for $a,b,c,d\in R$, if $a\sec(200^\circ)-c\tan(200^\circ)=d$ and $b\sec(200^\circ)+d\tan(200^\circ)=c$, then find the value of $\left(\dfrac{a^2+b^2+c^2+d^2}{bd-ac}\right)\sin20^\circ$.
\qn{7} If line $2x-y-1=0$ intersect the curve $ax^2+2hxy+by^2+2gx+2fy+c=0$ in two points $A$ and $B$, such that $\angle AOB=90^\circ$, then $\left|\dfrac{a+b+5c}{2g-f}\right|$ is equal to
\qn{8} The value of the expression $\dfrac1{\log_418}+\dfrac1{2\log_63+\log_62}+\dfrac5{\log_318}$ is
\end{multicols}
\end{document}
